\documentclass[tikz,border=10pt]{standalone}
\input{../../docs/tikz/dag_preamble.tex}

\begin{document}

% Agent visual validation loop:
% 1. Compile this file to PDF after every substantive DAG edit.
% 2. Inspect the PDF or a rendered PNG, not just the TeX source.
% 3. Increase spacing or reroute arrows until no box, legend, note, edge, or label overlaps.
%
% Intended lifecycle:
% - During paper intake, replace the scaffold below with every named result in the
%   paper and the dependency arrows between them.
% - After intake, treat this as a stable topology: update node styles/status/text
%   as proofs close, but avoid adding boxes or arrows unless the original named
%   result inventory or dependency map was actually incomplete.
%
% Arrow direction rule:
% Draw edges from prerequisite/source node to dependent/target node.  The
% `dag_arrow` and `dag_dashed_arrow` styles use an explicit end-arrow (`->`),
% so the arrowhead should always land on the node that consumes the dependency.
% Do not reverse the coordinate order just to make routing easier; use anchors
% or an orthogonal route through empty space instead.

\begin{tikzpicture}[
  x=1cm,
  y=1cm,
  every node/.append style={outer sep=4pt}
]

\dagPaperMetadata{Optimal Strategies in Ranked-Choice Voting}{Sanyukta Deshpande; Nikhil Garg; Sheldon H. Jacobson}{arXiv:2407.13661}{2024}{https://arxiv.org/abs/2407.13661}

% --- Legend: README status vocabulary plus formalized node-type styles. ---
\dagPaperLegendRightOfMetadata{
\node[dag_result, dag_template_legend] (legRes) at (0,0) {formalized\\result};
\node[dag_lemma, dag_template_legend] (legLem) at (4.2,0) {formalized\\lemma};
\node[dag_model, dag_template_legend] (legDef) at (8.4,0) {formalized\\definition};
\node[dag_caveat_legend] (legCav) at (12.6,0) {formalized\\with caveat};
\node[dag_partial, dag_template_legend] (legPart) at (0,-2.6) {partially\\formalized};
\node[dag_conditional, dag_template_legend] (legCond) at (4.2,-2.6) {conditional};
\node[dag_scaffold, dag_template_legend] (legScaf) at (8.4,-2.6) {scaffold};
\node[dag_unformalized, dag_template_legend] (legNot) at (12.6,-2.6) {not started /\\not formalized};
}
\daglegend{(legRes)(legLem)(legDef)(legCav)(legPart)(legCond)(legScaf)(legNot)}{Legend}

\begin{dagPaperBody}

\node[dag_model] (VotingModel) at (0,-3.8) {
  \textbf{RCV/STV Model} \\
  Ballots, traces, final orders, and win/loss sequences
};
\node[dag_model] (StructureConstraints) at (7.0,-3.8) {
  \textbf{Algorithms 1--2} \\
  Structure constraints, tally transfer, and direct-STV agreement
};
\node[dag_model] (SmartAllocation) at (14.0,-3.8) {
  \textbf{Algorithm 3} \\
  SmartAllocation slack filling, cost, and operation-count formulas
};
\node[dag_model] (CandidateReduction) at (21.0,-3.8) {
  \textbf{Algorithms 4--6} \\
  Reduced ballots, Strict-Support, and irrelevant-candidate removal
};

\node[dag_result] (Prop21) at (0,-7.6) {
  \textbf{Proposition 2.1} \\
  Tie-broken structures partition voter data and recover the STV order
};
\node[dag_result] (ThmB1) at (7.0,-7.6) {
  \textbf{Theorem B.1} \\
  The STV mechanism is a well-defined social choice order
};
\node[dag_lemma] (LemmaB2) at (14.0,-7.6) {
  \textbf{Lemma B.2} \\
  Round-winning candidates are never final election losers
};
\node[dag_result] (Prop33) at (21.0,-7.6) {
  \textbf{Proposition 3.3} \\
  Feasible sequence families satisfy the paper's win-count bound
};

\node[dag_result] (Thm31) at (14.0,-11.4) {
  \textbf{Theorem 3.1} \\
  SmartAllocation reaches a target structure optimally in polynomial time
};
\node[dag_lemma] (LemmaC1) at (7.0,-11.4) {
  \textbf{Lemma C.1} \\
  Initial elimination order is irrelevant for later transferred votes
};
\node[dag_result] (Thm32) at (21.0,-11.4) {
  \textbf{Theorem 3.2} \\
  Irrelevant candidates can be removed without changing later dynamics
};
\node[dag_model] (SequenceReduction) at (28.0,-7.6) {
  \textbf{Algorithm 7} \\
  Predict-Wins, Predict-Losses, and sequence-bound source formulas
};

\node[dag_result] (Prop34) at (21.0,-15.2) {
  \textbf{Proposition 3.4} \\
  Sequence search space can be shortened with quadratic profile runtime
};
\node[dag_model] (BenefitModel) at (0,-15.2) {
  \textbf{Definition 5.1} \\
  Benefit predicates, coalition benefit, and action model
};
\node[dag_result] (Prop53) at (7.0,-15.2) {
  \textbf{Proposition 5.3} \\
  Under perfect information, coalitions cannot all benefit from strategic voting
};
\node[dag_result] (Thm54) at (14.0,-15.2) {
  \textbf{Theorem 5.4} \\
  Optimal vote-addition strategies are characterized by Case (A)
};

\node[dag_result] (Prop55) at (7.0,-19.0) {
  \textbf{Proposition 5.5} \\
  Under uncertainty, coalitions may benefit ex ante but not all ex post
};
\node[dag_result] (Prop56) at (17.5,-19.0) {
  \textbf{Proposition 5.6} \\
  Selfish first-place additions are robust; other strategies may backfire
};
\draw[dag_arrow] (VotingModel.east) -- (StructureConstraints.west);
\draw[dag_arrow] (StructureConstraints.south west) -- (Prop21.north east);
\draw[dag_arrow] (StructureConstraints.south) -- (ThmB1.north east);
\draw[dag_arrow] (Prop21.east) -- (ThmB1.west);
\draw[dag_arrow] (ThmB1.east) -- (LemmaB2.west);
\draw[dag_arrow] (LemmaB2.east) -- (Prop33.west);
\draw[dag_arrow] (SmartAllocation.south east) to[out=-80,in=80] (Thm31.north east);
\draw[dag_arrow] (StructureConstraints.south east) to[out=-35,in=145] (Thm31.north west);
\draw[dag_arrow] (ThmB1.south) -- (LemmaC1.north);
\draw[dag_arrow] (CandidateReduction.south east) to[out=-80,in=80] (Thm32.north east);
\draw[dag_arrow] (LemmaC1.north east) .. controls +(3.0,1.5) and +(-3.0,1.5) .. (Thm32.north west);
\draw[dag_arrow] (SequenceReduction.south west) -- (Prop34.north east);
\draw[dag_arrow] (Prop33.south east) -- (Prop34.north west);
\draw[dag_arrow] (Thm32.south east) -- (Prop34.north west);
\draw[dag_arrow] (BenefitModel.east) -- (Prop53.west);
\draw[dag_arrow] (Thm31.south east) -- (Thm54.north west);
\draw[dag_arrow] (Prop53.south) -- (Prop55.north);
\draw[dag_arrow] (Thm54.south east) -- (Prop56.north west);
\draw[dag_arrow] (BenefitModel.south east) -- (Prop55.north west);
\end{dagPaperBody}
\end{tikzpicture}
\end{document}
